\begin{abstract}
This study investigates how server-side rendering (SSR) and client-side rendering (CSR) affect mobile client-side energy use within one comparable Next.js prototype. Although these rendering strategies are often discussed in terms of speed, interactivity, and search visibility, direct energy-oriented comparison in a controlled mobile setting remains limited.

To address this gap, one storefront application was implemented in two functionally comparable variants: an SSR route and a CSR route that shared the same interface, data model, and scenario structure. The final benchmark was executed on a Samsung Galaxy A53 under battery-powered wireless conditions, with 60 randomised runs across three dataset scenarios and two rendering conditions. Battery-derived energy was estimated by integrating current and voltage samples over time, while browser-side metrics such as LCP, FCP, TTFB, navigation response size, and a JavaScript heap sample were collected through the Chrome DevTools Protocol.

The corrected final dataset shows that SSR produced lower median total energy in all three scenarios, with the clearest support in the static and massive workloads. CSR, however, showed a smaller measured navigation response size in every scenario and showed lower paint-related timings in the dynamic case. The main contribution of this report is therefore not a universal winner, but a clearer account of the trade-off between server-shifted and client-shifted work when mobile energy efficiency is treated as a primary concern.

\end{abstract}
